\documentclass[10pt]{article}
\usepackage[margin=0.82in]{geometry}
\usepackage[T1]{fontenc}
\usepackage{lmodern}
\usepackage{amsmath,amssymb,amsthm}
\usepackage{graphicx}
\usepackage[hidelinks]{hyperref}

\newtheorem{theorem}{Theorem}
\newtheorem{corollary}[theorem]{Corollary}
\newcommand{\Iso}{\operatorname{Iso}_{\mathrm{lat}}}

\title{Compact Lattice-Isometry Groups Survive\\
the Benyamini Approximation}
\author{Automated mathematical discovery run \texttt{fa\_banach\_001}\\
Lane 6, \texttt{agent\_lane\_06}}
\date{9 August 2026}

\begin{document}
\maketitle

\noindent\textbf{Status.}
Partial affirmative result, likely valid, subject to human review.

\section{Source question}

Oikhberg and Tursi prove that, for every separable AM-space $X$ and every
$C>1$, there are a Benyamini space $X_C$ and a surjective lattice
isomorphism $A:X\to X_C$ satisfying
\[
 \|x\|\leq \|Ax\|\leq C\|x\| \qquad (x\in X).
\]
In Remark~2.4 they ask whether the construction can preserve the group of
lattice isometries, or even a subgroup \cite{source}.

\begin{center}
\includegraphics[width=0.67\textwidth]{figures/open_problem_crop.png}
\end{center}

Here ``preserve'' is made precise as the inclusion
$A G A^{-1}\subseteq\Iso(X_C)$. The proof below gives this for every subgroup
fixing a quasi-interior point and hence for every compact subgroup.

\newpage
\section{Result}

\begin{theorem}[Equivariant Benyamini approximation]\label{thm:fixed}
Let $X$ be a separable real AM-space, let $C>1$, and let
$G\leq\Iso(X)$. Suppose that $G$ fixes a quasi-interior point
$e\in X_+$. Then Proposition~2.3 of \cite{source} may be carried out so
that its lattice isomorphism $A:X\to X_C$ satisfies
\[
 \|x\|\leq\|Ax\|\leq C\|x\|,
 \qquad AGA^{-1}\subseteq\Iso(X_C).
\]
Conjugation by $A$ is also a topological group embedding for the strong
operator topologies.
\end{theorem}

\begin{corollary}\label{cor:compact}
Every strong-operator compact subgroup of $\Iso(X)$ can be preserved in the
sense of Theorem~\ref{thm:fixed}. In particular, every finite subgroup can be
preserved. If $\Iso(X)$ itself is compact, the entire original isometry group
survives as a subgroup of $\Iso(X_C)$.
\end{corollary}

\section{Proof}

\subsection*{The fixed-point criterion}

Use the canonical equivariant Kakutani representation. Namely, let
\[
 H=\{\varphi\in B_{X^*}:\varphi\text{ is a lattice homomorphism}\}
\]
with its weak-star topology, and put $(Jx)(\varphi)=\varphi(x)$. The set $H$
is compact, and the standard representation theorem for AM-spaces says that
$J:X\to C(H)$ is a lattice isometry. Its common zero is $0\in H$.

After scaling, assume $\|e\|\leq1$. If $\varphi\in H\setminus\{0\}$, then
$\varphi(e)>0$: otherwise $\varphi$ vanishes on the principal ideal generated
by $e$, hence on its dense closure $X$. Thus the function
$\psi=Je$ vanishes precisely at $0$.

For $n\geq1$, form the compact level sets
\[
 H_n=\{\varphi\in H:C^{-n}\leq\varphi(e)\leq C^{1-n}\}.
\]
Take disjoint copies $\widetilde H_n$, add the point at infinity when needed,
and let $q:\widetilde H\to H$ forget the copy index and send infinity to
$0$. This is the level construction used in the proof of Proposition~2.3.
Define, for $\widetilde\varphi_n\in\widetilde H_n$,
\[
 (Sx)(\widetilde\varphi_n)
   =\frac{C^{1-n}\varphi(x)}{\varphi(e)},
 \qquad (Sx)(\infty)=0.
\]
The scalar multiplier lies in $[1,C]$. Since $J$ is an isometry and every
point of $H\setminus\{0\}$ occurs in a level,
\[
 \|x\|\leq\|Sx\|_\infty\leq C\|x\|.
\]
Continuity at infinity follows as in the source proof: if the level indices
tend to infinity, then $\varphi(e)\to0$; quasi-interiority forces every
weak-star cluster point of the corresponding $\varphi$'s to be $0$, hence
$\varphi(x)\to0$, while the displayed multiplier remains bounded by $C$.

For $g\in G$, define the homeomorphism
\[
 \alpha_g:H\to H,\qquad \alpha_g(\varphi)=\varphi\circ g.
\]
Because $g$ is a lattice isometry and $ge=e$, this map preserves every
$H_n$. It therefore lifts to a homeomorphism $\widetilde\alpha_g$ of
$\widetilde H$, fixing infinity. Directly from the definition,
\begin{equation}\label{eq:equivariance}
 S(gx)=Sx\circ\widetilde\alpha_g \qquad (x\in X).
\end{equation}

Complete the source construction by identifying, within each level, points
on which every member of $S(X)$ has the same value. Equation
\eqref{eq:equivariance} shows that this equivalence relation is invariant
under every $\widetilde\alpha_g$. Hence these homeomorphisms descend to the
quotient compact pieces $K_n$. The source argument now gives a Benyamini
space $X_C\subset C_0(K)$ and a lattice isomorphism $A:X\to X_C$ with the
displayed distortion bounds. On $X_C$, the conjugate $AgA^{-1}$ is
composition with the descended homeomorphism, and therefore is a surjective
lattice isometry. This proves the algebraic assertion. Conjugation by the
fixed bounded isomorphisms $A$ and $A^{-1}$ is continuous in both directions
for point-norm convergence, proving the strong-operator statement.

\subsection*{Compact groups supply the fixed point}

Let $G\leq\Iso(X)$ be compact in the strong operator topology. Since $X$ is
separable, choose a quasi-interior point $e_0\in X_+$. With $m$ normalized
Haar measure on $G$, the Bochner integral
\[
 e=\int_G g e_0\,dm(g)
\]
exists, is positive, and is fixed by $G$. It remains to check that $e$ is
quasi-interior.

Let $0\neq f\in X^*_+$. Quasi-interiority of $e_0$ gives $f(e_0)>0$. The
function $g\mapsto f(ge_0)$ is continuous and nonnegative. If its Haar
integral were zero, it would vanish on all of $G$, contradicting its positive
value at the identity. Hence $f(e)>0$ for every nonzero $f\in X^*_+$. This
strict positivity characterizes quasi-interior points: otherwise the proper
closed ideal generated by $e$ would admit a nonzero positive functional on
the quotient that annihilates $e$. Thus $e$ is quasi-interior, and
Theorem~\ref{thm:fixed} applies.

\section{Scope and novelty}

The theorem supplies a faithful conjugate copy of $G$ in the new isometry
group. It does not assert that no additional lattice isometries appear. It
also leaves open arbitrary noncompact subgroups that have no fixed
quasi-interior point.

A bounded search on 9 August 2026 covered the run indexes, the parsed arXiv
corpus, exact-phrase searches for Remark~2.4, and close searches combining
Benyamini spaces, compact groups, and quasi-interior points. It found the
source and the authors' later $C_0(L)$ group-display paper \cite{followup},
which explicitly asks for a general AM-space extension, but no occurrence of
the equivariant criterion or the compact-subgroup corollary above. Novelty
confidence is moderate pending expert review.

\begin{thebibliography}{9}
\bibitem{source}
T. Oikhberg and M. A. Tursi,
\emph{Renorming AM-spaces}, Proc. Amer. Math. Soc. \textbf{150} (2022),
1127--1139; arXiv:2011.04140.

\bibitem{followup}
T. Oikhberg and M. A. Tursi,
\emph{Lattice renormings of $C_0(X)$ spaces}, Banach J. Math. Anal.
\textbf{18} (2024), article 76; arXiv:2405.11148.
\end{thebibliography}

\end{document}
